% Generated draft: requires mathematical and source review.
\documentclass[11pt]{article}
\usepackage[margin=25mm]{geometry}
\usepackage{fontspec}
\setmainfont{Latin Modern Roman}
\usepackage{amsmath,amssymb,mathtools,unicode-math}
\setmathfont{Latin Modern Math}
\usepackage{enumitem,needspace,xurl,hyperref}
\setlength{\parindent}{0pt}
\setlength{\parskip}{5pt}
\setlength{\emergencystretch}{3em}
\newcommand{\N}{\mathbb{N}}
\newcommand{\Z}{\mathbb{Z}}
\newcommand{\Q}{\mathbb{Q}}
\newcommand{\R}{\mathbb{R}}
\newcommand{\C}{\mathbb{C}}
\newcommand{\F}{\mathbb{F}}
\newcommand{\A}{\mathbb{A}}
\newcommand{\PP}{\mathbb{P}}
\newcommand{\E}{\mathbb{E}}
\newcommand{\eps}{\varepsilon}
\newcommand{\dd}{\,\mathrm{d}}
\begin{document}
\section*{UnsolvedMath ID: 30000220}
\textbf{Source problem code:} catalogue entry
\subsection{Definitions and mathematical statement}
Let \(R\) be a Dedekind domain and \(M\) an additive submonoid of a rational vector space \(V\), containing \(0\). Its monoid ring \(R[M]\) consists of finite sums \(\sum_{m\in M}a_mX^m\), with \(X^mX^{m'}=X^{m+m'}\). For integers \(c\geq2\) and \(t\geq0\), define
\[
 M^{1/c^t}=\{v\in V:c^tv\in M\}=c^{-t}M.
\]
Since \(M\subseteq c^{-t}M\), the inclusions induce \(R[M]\subseteq R[M^{1/c^t}]\). These are roots of monomial exponents; no roots of coefficients, assumption on the characteristic of \(R\), or restriction that \(c\) be a prime is inserted.

For a commutative ring \(B\), let \(SL_n(B)=\{A\in\operatorname{Mat}_n(B):\det A=1\}\), and let \(E_n(B)\) be the subgroup generated by \(e_{ij}(b)=I_n+bE_{ij}\), where \(i\ne j\) and \(b\in B\). Here \(E_{ij}\) is a matrix unit. Fix integers \(n\geq3\), \(c\geq2\) and \(A\in SL_n(R[M])\). The source reports that some finite \(t\in\N\) satisfies
\[
 A\in E_n(R[M^{1/c^t}]).
\]
The first question, identified in the original wording as a conjectural strengthening of Theorem 1.1, is whether one can always take \(t=1\). Explicitly, for every stated choice of \(R,M,n,c,A\), does there exist a finite factorization
\[
 A=\prod_{\ell=1}^{s}e_{i_\ell j_\ell}(b_\ell),
 \qquad b_\ell\in R[M^{1/c}],\quad i_\ell\ne j_\ell?
\]
The matrix \(A\) here is its image under the ring inclusion; no stabilization to larger matrices or alteration of \(A\) is allowed. The number of factors is not prescribed, and the parameter \(c\) is fixed before asking for the factorization.

The second, separate question asks whether the effective factorization procedure over a suitable root extension can be accelerated. It is not limited in the wording to a one-step extension. No algorithm, input encoding, cost model, asymptotic baseline, desired complexity bound or permitted operations on a general Dedekind domain are supplied. Thus an acceleration target cannot be made quantitative without further source data.

Theorem 1.1 itself is absent; the cleaned record supplies the hypotheses above and claims direct verification, but any further hypotheses in that theorem cannot be checked here. The literature assessment checked 2026-08-21 reports existence of a finite root exponent and an effective algorithm as established, while not resolving the one-step conjecture in the checked work. Its term ``Frobenius-root'' is retained as literature terminology, not used to impose positive characteristic on the displayed monoid construction. No independent proof or current status assessment is asserted.
\subsection{Short English statement}
The first question asks whether every stated special-linear matrix becomes elementary after adjoining only one level of monomial roots, rather than some unspecified finite number. The second asks for a faster effective factorization procedure. The ring, monoid, matrix size and fixed root parameter are retained; the algorithmic meaning of faster is not specified.
\subsection{Sources}
\begin{itemize}
\item Local UnsolvedMath registry, permanent ID 30000220, source code catalogue entry, ``Effective Factorization in Monoid Rings,'' including the reference to Theorem 1.1 and both original and expanded questions.
\item \emph{Kommutative Algebra (2006)}, Oberwolfach Reports, pp.\ 1073--1126, DOI 10.4171/owr/2005/19: \url{https://doi.org/10.4171/owr/2005/19}, also labelled ``original Oberwolfach Report 19/2005.'' The citation uses 2006, while the DOI and proposed-year metadata use 2005.
\item Supplied reference ``Gubeladze, Steinberg groups and algorithms'': \url{https://arxiv.org/abs/math/0601400}.
\item Literature triage checked 2026-08-21; registry record updated 2026-09-16. No external source has been visited or independently checked here.
\end{itemize}
\end{document}
